\section{Unitary evolution of a closed quantum system (9/9)}

\referto{\sakurai~2.1; \tong~3.3; \nielsen~2.1.8, 2.2.2, 4.2; \shankar~4.3}

In addition to Postulates \ref{postulate 1} and \ref{postulate 3:projective measurements} introduced earlier, we have two more postulates, one regarding how a quantum system evolves with time:
\begin{pos}\label{postulate 2}
    The time evolution of the state of a closed quantum system is described by the Schrodinger equation:
    \begin{align}\label{schrodinger eq}
    \imth\hbar\frac{\text{d}\dket{\psi}}{\text{d}t}=\hat H\dket{\psi}
    \end{align}
    where $\hbar=h/2\pi\approx1.05\times10^{-34}$ J$\cdot$s is a constant, and $\hat H$ is a fixed Hermitian operator known as the Hamiltonian of the closed system.
\end{pos}

Postulate \ref{postulate 2} states that the state vector associated with a closed quantum system evolves following the Schrodinger equation (\ref{schrodinger eq}). This postulate can be states in an equivalent manner:

\textbf{Postulate 3a} \emph{The evolution of a closed quantum system is described by a unitary transformation. That is, the state $\dket{\psi(t_2)}$ of the system at time $t_2$ is related to the state $\dket{\psi(t_1))}$ at time $t_1$ by a unitary operator $U$ which only depends on the times $t_1$ and $t_2$:}
\begin{align}
\dket{\psi(t_2)}=U(t_1,t_2)\dket{\psi(t_1)},~~~U(t_1,t_2)=e^{-\imth(t_2-t_1)\hat H/\hbar}
\end{align}
One can further generalize this to a time-dependent Hamiltonian $\hat H(t)$ where we have
\begin{align}
U(t_1,t_2)=\mathcal{T}e^{-\imth\int_{t_1}^{t_2}\hat H(t)\text{d}t/\hbar}
\end{align}
where $\mathcal{T}$ denotes time ordering in the integral, as explained in \sakurai~2.1.

In linear algebra, a unitary operator $\hat U$ is defined as any linear operator satisfying
\begin{align}
\hat U^\dagger\hat U=\hat 1.
\end{align}
which also implies $\hat U\hat U^\dagger=\hat 1$. Due to the spectrum decomposition theorem, a unitary operator can be diagonalized, and its eigenvalues all have modulus 1:
\begin{align}
\hat U\dket{\psi}=e^{\imth\theta}\dket{\psi}
\end{align}

\begin{table}[t]
\centering
\begin{tabular}{|c|c|c|}
\hline
Single-qubit gate&Symbol&Unitary matrix\\
\hline
Pauli-X (bit flip)&$X$&$\bpm0&1\\1&0\epm$\\
\hline
Pauli-Y&$Y$&$\bpm0&-\imth\\ \imth&0\epm$\\
\hline
Pauli-Z (phase flip)&$Z$&$\bpm1&0\\0&-1\epm$\\
\hline
Hadamard&$H$&$\frac1{\sqrt2}\bpm1&1\\1&-1\epm$\\
\hline
Phase&$S$&$\bpm1&0\\0&\imth\epm$\\
\hline
$\pi/8$&$T$&$\bpm1&0\\0&e^{\imth\pi/4}\epm$\\
\hline
\end{tabular}
\caption{Common single-qubit gates and associated unitary matrices.}\label{tab:single qubit gate}
\end{table}

Three examples of unitary operators in quantum mechanics are the following: (1) unitary time evolution operator $U(t_1,t_2)=e^{-\imth(t_2-t_1)\hat H/\hbar}$ for a closed quantum system discussed above; (2) a symmetry operator in a quantum system as we will discuss later in this course; (3) a quantum (logic) gate in quantum computing. For instance, in a two-dimensional closed quantum system, also known as a qubit, the common quantum gates are summarized in Table \ref{tab:single qubit gate}.

Starting from a complete set of orthonormal basis vectors $\{\dket{i}\}$ of Hilbert space $\mathbb{V}$, after the action of a unitary operator $\hat U$, we can obtain another complete set of orthonormal basis vectors $\{\dket{i}^\prime=\hat U\dket{i}\}$ of the same Hilbert space. In other words, each unitary operator implements a change of basis in the Hilbert space. Therefore, the unitary time evolution merely performs a rotation for vectors in the Hilbert space, since it only changes the state itself but not the inner product of any two states.  

In the case of unitary time evolution $U(t_1,t_2)$, the Hamiltonian $\hat H$ of a closed quantum system plays a crucial role. In its spectral decomposition:
\begin{align}
\hat H=\sum_aE_a\dket{a}\dbra{a}
\end{align}
the eigenkets $\dket{a}$ are called the \emph{energy eigenstates}, or \emph{stationary states}, and the eigenvalue $E_a$ is called the energy of state $\dket{a}$. In particular, the eigenstate with the lowest energy is known as the \emph{ground state}, while other states with higher energy are known as \emph{excited states}. According to Schrodinger equation (\ref{schrodinger eq}), a generic state vector
\begin{align}
\dket{\psi}=\sum_ac_a\dket{a}
\end{align}
evolves as
\begin{align}
\dket{\psi(t)}=e^{-\imth t\hat H/\hbar}\dket{\psi}=\sum_ac_ae^{-\imth E_at/\hbar}\dket{a}
\end{align}
and the different oscillation frequency $\omega_a=E_a/\hbar$ in different components of state $\dket{\psi(t)}$ can result in oscillations of physical observables. Two such examples, i.e. (1) spin precession; and (2) neutrino oscillation are discussed in detail in \sakurai~2.1, and will be a part of HW3. In particular, a physical observable you will encounter a lot in the future is the \emph{transition probability}:
\begin{align}
P(i\rightarrow f)\equiv|\expval{f|e^{-\imth\hat Ht/\hbar}|i}|^2
\end{align}
which measures the probability that an initial state $\dket{i}$ evolves into a final state $\dket{f}$ after time $t$. This quantity is directly related to the outgoing beam intensity in scattering experiments.

The transition probability is the modulus square of the \emph{propagator}
\begin{align}
    G(i,f:t)\equiv\expval{f|e^{-\imth\hat Ht/\hbar}|i}
\end{align}
which plays very important conceptual and computational roles in quantum mechanics. We will learn more about the propagator in HW3. 

\subsection{Schrodinger vs. Heisenberg pictures}

Regarding the time dependence of physical observables, such as the expectation value of a Hermitian operator $\hat O$ at time $t$, there are two equivalent ways to approach this problem. 

In the Schrodinger picture, where the state vector $\dket{\psi}$ evolves with time while the operator $\hat O$ remains a constant:
\begin{align}
    \dket{\psi_S(t)}=\hat U(t,0)\dket{\psi_0},~~~\hat O_S(t)=\hat O
\end{align}
In the Heisenberg picture, on the other hand, the state vector remains a constant while the operator evolves with time:
\begin{align}
    \dket{\psi_H(t)}=\dket{\psi_0},~~~\hat O_H(t)=\hat U^\dagger(t,0)\hat O\hat U(t,0)
\end{align}
It is straightforward to show that both pictures yield the same physical observable
\begin{align}
    \expval{O(t)}=\expval{\psi_S(t)|\hat O_S(t)|\psi_S(t)}=\expval{\psi_H(t)|\hat O_H(t)|\psi_H(t)}=\expval{\psi_0|\hat U^\dagger(t,0)\hat O\hat U(t,0)|\psi_0}
\end{align}

In general, it is more convenient to use the Heisenberg picture to compute expectation values of an observable, while the Schrodinger picture is more convenient to extract time dependence of the state vector and its related properties, such as entanglement entropy.

\subsection{Quantum Zeno effect}

Consider a quantum system under repeated measurements of a physical observable $\hat O$ with the following spectral decomposition:
\begin{align}
    \hat O=\sum_q q~\dket{q}\dbra{q}
\end{align}
Without loss of generality, let's assume the initial state of the system is $\hat O$ eigenstate $\dket{0}$ with eigenvalue $q_0$. 

Now consider a generic Hamiltonian
\begin{align}
    \hat H=H_{0,0}\dket0\dbra0+\sum_{q\neq0}(H_{q,0}\dket{q}\dbra0+h.c.)+\sum_{p,q\neq0}H_{p,q}\dket{p}\dbra{q}
\end{align}
that governs the time evolution of the system. We used $h.c.$ as a shorthand notation for ``hermitian conjugate''. 

Starting from the initial state $\dket0$, after a very short time $\Delta t$, the state evolves into
\begin{align}
    \dket{\psi(t)}=e^{-\imth \hat Ht/\hbar}\dket0=(1-\frac\imth\hbar H_{0,0}\Delta t)\dket0-\frac\imth\hbar\Delta t\sum_q H_{q,0}\dket{q}+O(\Delta t^2)
\end{align}
Therefore at time $\Delta t$, the probability of measurement outcome $\hat O=q_0$ is given by
\begin{align}
    |\expval{0|\psi(t)}|^2=1-\alpha(\Delta t)^2+O(\Delta t^3),~~~\alpha>0
\end{align}
where the constant $\alpha$ depends on the Hamiltonian $\hat H$. 

One can further show that if we repeat this 2-step procedure of small time evolution for $\Delta t=\frac TN$ followed by a measurement of operator $\hat O$, after $N$ repetitions, the probability for the state to come back to itself approaches 1 in the large $N$ limit:
\begin{align}
    \text{Pr}(\dket{0}\rightarrow\dket0)\geq\lim_{N\rightarrow\infty}\big(|\expval{0|\psi(\frac TN)}|^2\big)^N=\lim_{N\rightarrow+\infty}e^{-\alpha T^2/N}=1
\end{align}
This is known as the quantum Zeno effect. 